% !PRINTED 104
\weberpage{104}
To determine the constant $A$ we employ, as we shall often do in what follows, the device of differentiating (1), setting $u = 0$, and then using formula (5), § 23, on both sides:
\weq{5}{\vartheta'_{11} = \pi\,\vartheta_{00}\,\vartheta_{10}\,\vartheta_{01}}
It follows that
\[
A^2 = e^{\frac{\pi i}{2}}, \quad A = e^{\frac{\pi i}{4}};
\]
that the positive sign stands with $A$ follows from one of the formulae (3), (4), according to the closing remark of § 25, if one lets $\omega$ become infinite.

Accordingly one obtains
\weq{6}{\begin{aligned}
\vartheta_{11}(u, \omega+1) &= e^{\frac{\pi i}{4}}\vartheta_{11}(u) \\
\vartheta_{10}(u, \omega+1) &= e^{\frac{\pi i}{4}}\vartheta_{10}(u) \\
\vartheta_{01}(u, \omega+1) &= \vartheta_{00}(u) \\
\vartheta_{00}(u, \omega+1) &= \vartheta_{01}(u).
\end{aligned}}
II. $\begin{pmatrix} 0, & 1 \\ -1, & 0 \end{pmatrix}$ or $\left(\omega, \dfrac{-1}{\omega}\right)$.

Again we have, by § 27, (12),
\weq{7}{e^{-\frac{\pi i u^2}{\omega}}\vartheta_{11}\left(\frac{u}{\omega}, -\frac{1}{\omega}\right) = A\vartheta_{11}(u, \omega),}
and, on increasing $u$ by $\dfrac{1}{2}$, $\dfrac{\omega}{2}$, $\dfrac{1+\omega}{2}$,
\weq{8}{e^{-\frac{\pi i u^2}{\omega}}\vartheta_{01}\left(\frac{u}{\omega}, -\frac{1}{\omega}\right) = iA\vartheta_{10}(u, \omega),}
\weq{9}{e^{-\frac{\pi i u^2}{\omega}}\vartheta_{10}\left(\frac{u}{\omega}, -\frac{1}{\omega}\right) = iA\vartheta_{01}(u, \omega),}
\weq{10}{e^{-\frac{\pi i u^2}{\omega}}\vartheta_{00}\left(\frac{u}{\omega}, -\frac{1}{\omega}\right) = iA\vartheta_{00}(u, \omega),}
from which one obtains, as above:
\[
A = \pm i\sqrt{-i\omega}.
\]
From $u = 0$ and a purely imaginary $\omega$ one concludes that, if $\sqrt{-i\omega}$ is taken so that \gesperrt{the real part is positive, the lower sign must stand}, and there results therefore:
%CONT>>
